 \documentclass[11pt]{memoir}

% based on kieran healy's memoir modifications
\usepackage{mako-mem}
\chapterstyle{article-2}
\pagestyle{mako-mem}

% \usepackage{ucs}
% \usepackage[utf8x]{inputenc}

%\usepackage{kpfonts}
%\usepackage[bitstream-charter]{mathdesign}
\usepackage{fbb}
\usepackage[T1]{fontenc}
%\usepackage{textcomp}

%\renewcommand{\rmdefault}{ugm}
%\renewcommand{\sfdefault}{phv}

% Packages for making a landscape table (?)
\usepackage[table,dvipsnames]{xcolor}
\usepackage{multirow, makecell}
\usepackage{pdflscape}
\usepackage{afterpage}

\usepackage[letterpaper,left=1.125in,right=1.125in,top=1.25in,bottom=1.25in]{geometry}

% packages i use in essentially every document
\usepackage{graphicx}
\usepackage{enumerate}
\usepackage{stringstrings}
\usepackage{hwemoji}
\DeclareUnicodeCharacter{FE0F}{}

% date auto-incrementation
\usepackage{datenumber}
\setdate{2026}{8}{21}
\def\datedate{\datedayname,\space\datemonthname~\thedateday}
\def\datedateshort{\substring{\datemonthname}{1}{3}~\thedateday}

\newcommand{\adddays}[1]{%
    \addtocounter{datenumber}{#1}%
    \setdatebynumber{\thedatenumber}%
}

% Setup list environments
\usepackage{enumitem}
\setlist[description]{
  topsep=0pt,
  before=\vspace{0pt},
  after=\vspace{0pt},
  itemsep=2pt,
  labelsep=0pt
}

\setlist[itemize]{
    noitemsep,
    leftmargin=1em,
    topsep=0pt}

% packages i use in many documents but leave off by default
% \usepackage{amsmath, amsthm, amssymb}
% \usepackage{dcolumn}
% \usepackage{endfloat}

% Set paragraph indents and spacing
\setlength{\parindent}{0pt}
\setlength{\parskip}{.5\baselineskip}
%\usepackage[document]{ragged2e}

% adjust section title formatting
\usepackage{titlesec}
\titlespacing\section{0pt}{8pt plus 2pt minus 2pt}{-6pt}
\titlespacing\subsection{0pt}{8pt plus 2pt minus 2pt}{-6pt}
\titlespacing\subsubsection{0pt}{8pt plus 2pt minus 2pt}{10pt}

% allows full, in-line citations
\usepackage{bibentry}

% add bibliographic stuff
\usepackage[round, numbers]{natbib} \def\citepos#1{\citeauthor{#1}'s (\citeyear{#1})} \def\citespos#1{\citeauthor{#1}' (\citeyear{#1})}
\renewcommand{\bibnumfmt}[1]{}

% define colors from http://www.colorado.edu/brand/visual-identity/typography-color
%\usepackage[usenames,dvipsnames]{color}
\definecolor{CUGold}{RGB}{207,184,124}
\definecolor{CUDarkGray}{RGB}{86,90,92}
\definecolor{CULightGray}{RGB}{162,164,163}

% customize URLs
\usepackage[hyphens]{url}
%\usepackage{breakurl}
\usepackage[breaklinks, bookmarks, bookmarksopen]{hyperref}
\hypersetup{
    colorlinks=true,
    linkcolor=Blue,
    citecolor=Black,
    filecolor=Blue,
    urlcolor=Blue,
    unicode=true,
    breaklinks=true}

% create a "reading list" environment to format the items
\newenvironment{readinglist}{
\begin{list}{}{\leftmargin=0pt \itemindent=0em}
  \setlength{\itemsep}{8pt}
  \setlength{\parskip}{0em}
  \setlength{\parindent}{8em}}
{\end{list}}

% Course/Instructor metadata -- alter as needed
\def\coursename{Web Data Science}
\def\courselisting{INFO 4617}
\def\classroom{CASE W262}
% \def\zoomurl{https://cuboulder.zoom.us/j/91984080043}
% Canvas course URL (Fall 2026)
\def\canvasurl{https://canvas.colorado.edu/courses/143304}
\def\githuburl{https://github.com/cuinfoscience/INFO4617-Fall2026}
\def\bookurl{https://cuinfoscience.github.io/Web-Data-Science-Book/}
\def\bookgithuburl{https://github.com/cuinfoscience/Web-Data-Science-Book}
\def\meetingdays{Mondays, Wednesdays, and Fridays}
\def\starttime{12:20 pm}
\def\endtime{1:10 pm}
\def\meetingtimes{\starttime~-- \endtime}
\def\semester{Fall 2026}

\def\instructorfirstname{Brian}
\def\instructorlastname{Keegan}
\def\instructorfullname{\instructorfirstname~\instructorlastname, Ph.D.}
\def\instructortitle{Associate Professor, Information Science}
\def\instructoroffice{\href{https://www.colorado.edu/map/?id=336\#!m/193981}{INFO 129}}
\def\instructoremail{brian.keegan@colorado.edu}
\def\instructorwebsite{http://www.brianckeegan.com}
\def\instructorofficehours{Wednesdays, 10am-12pm}

\begin{document}

\nobibliography*

%\baselineskip 14.2pt

\title{
    \textbf{\huge{\coursename}}\\
    \vspace{5pt} \normalsize{\courselisting}; \semester
}

\author{
    \meetingdays; \meetingtimes\\
    % Zoom room: \href{\myzoomurl}{\myzoomurl}\\
    Classroom: \classroom
}

\date{
    \normalsize{
        \href{\instructorwebsite}{\textbf{\instructorfullname}}\\
        \instructortitle\\
        E-mail: \href{mailto:\instructoremail}{\instructoremail}\\
        Office: \instructoroffice\\
        Office hours: \instructorofficehours
    }
}

\maketitle

%%%%%%%%%%%%%%%%%%%%%%
%% Acknowledgements %%
%%%%%%%%%%%%%%%%%%%%%%
% This syllabus template was made in LaTeX by Brian Keegan and is distributed as Free Software under the GNU GPL v3. It was built using style templates created by Aaron Shaw, Benjamin Mako Hill, and Kieran Healy.

% Original description: https://docs.google.com/document/d/1l7Qaop6Lmy7wkP_Jdb5JkeXNkrWVAMdZgQgj5uyvN8Q/edit#

\section{Course Description}

The internet makes many kinds of information easy to access. The ability to retrieve, parse, and analyze this information is a valuable skill for data scientists. This course provides an overview of computational tools and practices for transforming web documents and APIs into data for common research designs.

The course is organized around an open-source textbook, the \href{\bookurl}{\textit{Web Data Science}} book, which we will read chapter-by-chapter and---crucially---improve together. Each week you will implement the concepts in companion Jupyter notebooks and then propose and evaluate revisions to the textbook itself through GitHub. By the end of the semester you will have built a portfolio of web-data skills, contributed to a living technical resource, and completed an original data-collection project.

\subsection{Learning objectives}

\begin{itemize}[itemsep=0em]
    \item Understand the history, infrastructure, legal and ethical contours of web data access
    \item Navigate and parse common web data formats like XML and JSON
    \item Retrieve and automate data extraction from HTML and PDF documents
    \item Access popular APIs to collect data for common research designs
    \item Use version control and open-source workflows to collaborate on a shared codebase
    \item Document, communicate, and critically evaluate data collection code and technical writing
\end{itemize}

\subsection{Course Design}
Class will meet three times per week on \meetingdays~from \starttime~to \endtime~in \classroom. The week follows a consistent rhythm:
\begin{description}[itemsep=2pt,labelsep=0.5em,leftmargin=1.5em]
    \item[Monday] introduces a new concept and its companion notebook through lecture and live coding.
    \item[Wednesday] is a hands-on Notebook Lab: you work through and share your implementation of the notebook exercises.
    \item[Friday] is a Textbook Revision workshop: you propose revisions to the Web Data Science book as GitHub pull requests and review your classmates' proposals.
\end{description}
Student performance is evaluated through participation in the Wednesday labs, the Friday textbook revisions, and a Final Project (see \textit{Evaluation} below). There is no midterm or final exam.

The class is organized into four modules that track the textbook: (1) \textit{Foundations} introduces the ethical and legal considerations of web data access, the ``post-API'' landscape, common file formats, and web protocols; (2) \textit{Documents} covers methods and tools for retrieving data from web pages and PDFs; (3) \textit{APIs} provides an overview of tools for accessing popular web APIs; and (4) \textit{Practice} covers automation, research design, and the Final Project. See Table~\ref{tab:course_outline} for the full outline.

\subsection{Prerequisites}
There are no formal pre-requisites for the class, but prior programming experience in Python is strongly recommended. If you have questions about these prerequisites, please \href{mailto:\instructoremail}{email the instructor}.

\subsection{Course Website and Materials}
Readings, notebooks, assignments, and grades are coordinated through Canvas:
\vspace{-8pt}
    \begin{center}
    \href{\canvasurl}{\canvasurl}
    \end{center}
\vspace{-8pt}

There is no textbook to purchase. The required text is the open-source \textit{Web Data Science} book, which we read and revise together. You can access the rendered version of the book here:
\vspace{-8pt}
    \begin{center}
    \href{\bookurl}{\bookurl}
    \end{center}
\vspace{-8pt}

The source code you will be editing, can be accessed here:
\vspace{-8pt}
    \begin{center}
    \href{\bookurl}{\bookgithuburl}
    \end{center}
\vspace{-8pt}

This PDF version of the syllabus is a work-in-progress and will be revised as we proceed through the semester. Any revised requirements will be posted as announcements and an updated course schedule to Canvas. The instructor reserves the right to make changes to the course's schedule, evaluation criteria, policies, \textit{etc.} through announcements in class and on Canvas. If there are any discrepancies between the syllabus and Canvas, defer to the latest instructions on Canvas. Students should \href{mailto:\instructoremail}{email the instructor} if there are any questions.

\subsection{Computing}
Students will use programming languages for data retrieval, analysis, and visualization. \href{http://jupyter.org/}{Jupyter notebooks} written in Python 3 will be used for all in-class examples and assignments. The \href{https://www.anaconda.com/}{Anaconda distribution} of Python 3.12 (or above) is \textit{strongly} recommended to provide these programs and libraries. Over the semester we will use libraries including \href{https://requests.readthedocs.io/}{\texttt{requests}}, \href{https://www.crummy.com/software/BeautifulSoup/bs4/doc/}{\texttt{BeautifulSoup}}, \href{https://selenium-python.readthedocs.io/}{\texttt{selenium}}, \href{https://pypdf.readthedocs.io/}{\texttt{pypdf}}, and \href{https://pandas.pydata.org/}{\texttt{pandas}}. You will also need \href{https://git-scm.com/}{Git} and a free \href{https://github.com/}{GitHub} account; we will set these up together in the first week. Lectures include exercises and presentations with the expectation that students participate with their own laptop computers. If you cannot bring a laptop to class, please email the instructor to work out an alternative arrangement.

\subsection{Evaluation}
Students are evaluated through three mechanisms. The class has no midterm or final exam.

    \begin{description}[itemsep=0pt,labelsep=0pt,leftmargin=1em]

        \item[Notebook Labs]~(30\%). The Wednesday labs are where you build fluency with retrieving, reshaping, and analyzing web data. Each week you work through and share your implementation of the companion notebook. Labs are graded on participation and completion rather than perfection, and are meant to be collaborative. Your lowest two lab scores are dropped to accommodate illness, travel, and the occasional bad week.

        \item[Textbook Revisions]~(15\%). The Friday workshops are dedicated to improving the open-source \textit{Web Data Science} book. Each week you propose revisions---clarifications, corrections, new examples, additional exercises, or better explanations---as pull requests to the book's GitHub repository, and you review your classmates' pull requests. Credit reflects both the substance of the revisions you propose and the quality of the peer reviews you provide. A running target for contributions and reviews across the semester will be posted to Canvas. In the absence of an approved excuse, late submissions will be docked 1\% of their maximum value for every day elapsed since the deadline.

        \item[Attendance]~(15\%). Attendance in class is required. We will be covering technical methods that are cumulative and require sustained effort. If there are personal, professional, medical, or other circumstances that will prevent your attendance, please notify me via e-mail and we can develop an accommodation plan together. There will be no opportunities to make up missed attendance. If you need to miss multiple classes, please make an effort to come to office hours so we can check in about course material and progress.

        \item[Final Project]~(40\%). The Final Project is a portfolio piece that matches a web data source with a research design. It is completed individually (or in pairs with the instructor's permission) and has three graded components: a \textbf{proposal} (due Friday, October 23), an in-class \textbf{presentation} during the final week, and a \textbf{project repository and write-up} (due Monday, December 7). Further details will be collaboratively developed and detailed later in the course. In the absence of an approved excuse, late submissions will be docked 1\% of their maximum value for every day elapsed since the deadline.

    \end{description}

Note that the Wednesday labs and Friday revision workshops are themselves graded on in-class participation. Missing one of those sessions therefore affects your grade twice: once through the Attendance component above, and again through that week's lab or revision score.

\section{Course Policies}

\subsection{In-Class Confidentiality}
The success of this class depends on students feeling comfortable sharing questions, ideas, concerns, and confusions about assignments, work-in-progress, and their personal experiences. Students may read, comment, and run classmates' writing, code, and other class-related content for use within this class. However, students may not use, run, copy, perform, display, distribute, modify, translate, or create derivative works of another student's work outside of this class without that student's expressed written consent or formal license. Students may not create any audio, video, or other records of lectures without the instructor's permission nor may students share comments made in class attributable to another person's identity without permission.

Note that our textbook contributions are different: pull requests to the \textit{Web Data Science} book are made to a \textit{public} open-source repository under an open license, and your accepted contributions will be visible to anyone. If you prefer to contribute under a pseudonym or have concerns about public attribution, email the instructor and we will work out an alternative.

\subsection{Acceptable use of AI}
Generative AI tools (\textit{e.g.}, ChatGPT, Claude, GitHub Copilot) are part of contemporary data science practice. You can use them in this course as collaborators subject to four expectations apply. First, \textbf{disclosure}: when AI tools shape anything you submit (a notebook, a pull request, \textit{etc}.), note where and how you used them (a brief comment or commit message is sufficient). Second, \textbf{understanding}: you are responsible for every line you submit and must be able to explain how your code works and why it is correct; the instructor may ask you to do so. Third, \textbf{integrity}: presenting AI-generated work as another person's, or using it to misrepresent data, sources, or effort, violates the Honor Code. When in doubt, disclose and ask. Fourth, \textbf{reflect}: critically examine the limitations of AI-generated work in terms of its accuracy, accessibility, and impacts on people.

\subsection{Classroom Behavior}
Students and faculty are responsible for maintaining an appropriate learning environment in all instructional settings, whether in person, remote, or online. Failure to adhere to such behavioral standards may be subject to discipline. Professional courtesy and sensitivity are especially important with respect to individuals and topics dealing with race, color, national origin, sex, pregnancy, age, disability, creed, religion, sexual orientation, gender identity, gender expression, veteran status, marital status, political affiliation, or political philosophy. For more information, see the policies on \href{http://www.colorado.edu/policies/student-classroom-and-course-related-behavior}{class behavior}, the \href{http://www.colorado.edu/osc/#student_code}{Student Code of Conduct}, and the \href{https://www.colorado.edu/oiec/}{Office of Institutional Equity and Compliance}.

\subsection{Mental Health and Wellness}
The University of Colorado Boulder is committed to supporting students’ mental health and overall well-being. If personal, academic, or emotional challenges are affecting your well-being or success, \href{https://www.colorado.edu/counseling/}{Counseling and Psychiatric Services} (CAPS), is here to help. CAPS offers counseling, referrals, psychiatric care, crisis support, and much more. Visit CAPS in the C4C, or call (303) 492-2277, 24/7.

\subsection{Accommodations for Disabilities}
If you qualify for accommodations because of a disability, please submit your accommodation letter from Disability Services to your faculty member in a timely manner so that your needs can be addressed.  Disability Services determines accommodations based on documented disabilities in the academic environment.  Information on requesting accommodations is located on the \href{www.colorado.edu/disabilityservices/students}{Disability Services website}. Contact Disability Services at 303-492-8671 or \href{mailto:dsinfo@colorado.edu}{dsinfo@colorado.edu} for further assistance. If you have a temporary medical condition, see \href{http://www.colorado.edu/disabilityservices/students/temporary-medical-conditions}{Temporary Medical Conditions} on the Disability Services website.

If you have a temporary illness, injury or required medical isolation for which you require adjustment, please contact the instructor as soon as possible to develop an accommodation plan. You are not required to state the nature of the illness nor do you need to share a doctor's note.

\subsection{Student Names and Pronouns}
CU Boulder recognizes that students' legal information does not always align with how they identify. If you wish to have a name other than your legal name appear on your instructors’ class rosters and in Canvas, or if you wish to choose pronouns to appear on your instructors’ class rosters and in Canvas, visit the \href{https://www.colorado.edu/registrar/students/records/info/preferred}{Registrar’s website} for instructions on how to change your personal information in university systems.

\subsection{Honor Code}
All students enrolled in a University of Colorado Boulder course are responsible for knowing and adhering to the \href{http://www.colorado.edu/policies/academic-integrity-policy}{Honor Code}. Violations of the Honor Code may include but are not limited to: plagiarism (including use of paper writing services or technology [such as essay bots]), cheating, fabrication, lying, bribery, threat, unauthorized access to academic materials, clicker fraud, submitting the same or similar work in more than one course without permission from all course instructors involved, and aiding academic dishonesty. Understanding the course's syllabus is a vital part of adhering to the Honor Code. All incidents of academic misconduct will be reported to the Student Conduct \& Conflict Resolution (\href{mailto:studentconduct@colorado.edu}{studentconduct@colorado.edu}; 303-492-5550). Students found responsible for violating the Honor Code will be assigned resolution outcomes from Student Conduct \& Conflict Resolution and will be subject to academic sanctions from the faculty member. Visit \href{https://www.colorado.edu/sccr/students/honor-code-and-student-code-conduct}{Honor Code} website for more information on the academic integrity policy. 

\subsection{Harassment and Discrimination}
CU Boulder is committed to fostering a productive and welcoming learning, working, and living environment. University policy prohibits \href{https://www.colorado.edu/oiec/policies/protected-class-nondiscrimination-policy/protected-class-definitions}{protected class} discrimination and harassment, sexual misconduct (harassment, exploitation, and assault), intimate partner abuse (dating or domestic violence), stalking, and related retaliation by or against members of our community on- or off-campus. Denial of an approved accommodation for disability, religious observance, or pregnancy or pregnancy related medical conditions may be discriminatory.

The \href{https://www.colorado.edu/oiec/}{Office of Institutional Equity and Compliance} (OIEC) addresses these concerns, and individuals who have been subjected to misconduct can contact OIEC at 303-492-2127 or email \href{mailto:OIEC@colorado.edu}{OIEC@colorado.edu}. Information about university policies, \href{https://www.colorado.edu/oiec/reporting-resolutions/making-report}{OIEC reporting options}, and \href{https://www.colorado.edu/oiec/support-resources}{OIEC support resources} including confidential services can be found on the \href{https://www.colorado.edu/oiec/}{OIEC website}.

Faculty and graduate instructors are required to inform OIEC when someone discloses misconduct regardless of when or where it occurred. This is to ensure that those impacted receive outreach from OIEC about resolution options and support resources. To learn more about reporting and support options for a variety of concerns, visit the \href{https://www.colorado.edu/dontignoreit/}{Don’t Ignore It} page.

\subsection{Religious Observance}
Instructional faculty members must make every reasonable effort to accommodate all students who have conflicts with scheduled exams, assignment deadlines or required attendance due to a religious observance. Whenever possible, students must notify the instructional faculty member at least two weeks in advance of the expected exam, assignment deadline, or attendance conflict to request an accommodation for religious observance. If the start date of the course is less than two weeks before the date of requested accommodation, students must notify the instructional faculty member on the start date of the course. See the \href{https://www.colorado.edu/compliance/student-academic-accommodations-religious-observances-policy
}{Student Academic Accommodations for Religious Observances Policy} for more information.  If you need to miss class for a religious observance, please \href{mailto:brian.keegan@colorado.edu}{email the instructor} as soon as possible to make the appropriate accommodations.

\subsection{CMCI Diversity, Equity, and Inclusion}
CMCI strives to be a community whose excellence depends on diversity, equity, and inclusion. We aim to understand and challenge systems of privilege and disadvantage in higher education, such as those based on class, race, ethnicity, gender, sexuality, and dis/ability. We seek to reach across social and political divides and to make space for voices historically underrepresented in higher education and marginalized in society. In other words, diversity is not just a future reality for which we try to prepare students. It is a priority we want to put into practice here, now, and together, in order to foster places of learning where all members can thrive. Please contact the CMCI diversity team (\href{mailto:lisa.flores@colorado.edu}{email Lisa Flores} or visit the \href{https://www.colorado.edu/cmci/about-college/diversity-equity-and-inclusion/our-team}{CMCI Diversity, Inclusion, and Equity Staff} page) if: you need support or other resources but don't know where to turn, any aspect of your educational experience with CMCI does not reflect the commitment expressed here, or if you want to share a positive instance of this commitment in action, you have any questions, concerns, or ideas related to diversity.

\section{Acknowledgments}
% The design and format of this course borrows from other courses.
This syllabus was typeset in \LaTeX~using \href{http://www.overleaf.com}{Overleaf} with the \href{http://www.tug.dk/FontCatalogue/fbb/}{fbb/Bembo} font and is derived from the \texttt{memoir} styles adapted by \href{https://github.com/kjhealy/latex-custom-kjh}{Kieran Healy} and \href{http://projects.mako.cc/source/?p=latex_mako;a=summary}{Benjamin `Mako' Hill}.

\section{Other resources}
\begin{itemize}
    \item \bibentry{broucke_PracticalWebScraping_2018}
    \item \bibentry{bradshaw_ScrapingJournalists_2012}
    \item \bibentry{chapagain_web_2019}
    \item \bibentry{heydt_python_2018}
    \item \bibentry{lawson_WebScrapingPython_2015}
    \item \bibentry{mitchell_WebScrapingPython_2018}
    \item \bibentry{munzert_AutomatedDataCollection_2014}
    \item \bibentry{rogers_DoingDigitalMethods_2019}
    \item \bibentry{russell_mining_2019}
    \item \bibentry{salganik_BitBitSocial_2019}
\end{itemize}

%%%%%%%%%%%%%%%%%%%%%%
%%%%%%%%%%%%%%%%%%%%%%
%%% COURSE OUTLINE %%%
%%%%%%%%%%%%%%%%%%%%%%
%%%%%%%%%%%%%%%%%%%%%%
\clearpage
\section{Course Outline}

This is an outline of modules and topics, mapped to the chapters of the \href{\bookurl}{\textit{Web Data Science}} book. The schedule will evolve; consult Canvas for the most up-to-date information.

\begin{table}[ht]
%\Large
\centering
    \begin{tabular}{cccc}
        \toprule[.15em]
        \textbf{Module} & \textbf{Week} & \textbf{Dates} & \textbf{Topics} \\
        \cmidrule[.1em](lr){1-4}

        \multirow{5}{*}[0pt]{\textit{Foundations}}
            & 1 & Aug.~21 & Introduction \& setup \\
            & 2 & Aug.~24--28 & Ethics, law \& responsible collection \\
            & 3 & Aug.~31--Sep.~4 & The post-API age \\
            & 4 & Sep.~9--11 & Data formats: XML \& JSON \\
            & 5 & Sep.~14--18 & Web architecture \& protocols \\ \cmidrule[.1em](lr){1-4}

        \multirow{4}{*}[0pt]{\textit{Documents}}
            & 6 & Sep.~21--25 & Parsing static web pages \\
            & 7 & Sep.~28--Oct.~2 & Archived web pages \\
            & 8 & Oct.~5--9 & Dynamic web pages \\
            & 9 & Oct.~12--16 & PDFs \\ \cmidrule[.1em](lr){1-4}

        \multirow{4}{*}[0pt]{\textit{APIs}}
            & 10 & Oct.~19--23 & Wikipedia APIs \\
            & 11 & Oct.~26--30 & Government data APIs \\
            & 12 & Nov.~2--6 & Social media APIs \\
            & 13 & Nov.~9--13 & AI APIs \\ \cmidrule[.1em](lr){1-4}

        \multirow{3}{*}[0pt]{\textit{Practice}}
            & 14 & Nov.~16--20 & Automating data collection \\
            & 15 & Nov.~23--27 & \textbf{No class: Fall Break} \\
            & 16 & Nov.~30--Dec.~4 & Final presentations \\
            \cmidrule[.1em](lr){1-4}
            & \multicolumn{3}{c}{\textbf{Final Project due Monday, December 7}} \\

        \bottomrule[.15em]
    \end{tabular}
    \caption{Course outline by week. Labor Day (Monday, Sep.~7) and Fall Break (Nov.~23--27) have no class; the last day of class (Dec.~4) follows a Monday schedule.}
    \label{tab:course_outline}
\end{table}

\clearpage

\section{\textbf{Course Schedule}}

The schedule will evolve throughout the semester, so please consult the schedule online at Canvas for the most up-to-date information. Unless noted otherwise, each week follows the same rhythm: \textbf{Monday} introduces the concept and notebook, \textbf{Wednesday} is a hands-on notebook lab, and \textbf{Friday} is a textbook-revision workshop.

\section{Week 1 -- Foundations: Introduction \& setup}
\textcolor{CUGold}{\textbf{Friday, August 21}}\\
Course overview and motivation; setting up the Python computing environment, Git, and GitHub; how our textbook-revision workflow will work. \textit{Because the semester begins on a Thursday, our first meeting is this Friday only.} Textbook: \href{\bookurl/blob/main/ch-01-introduction.qmd}{Chapter 1 (Introduction)}.

\section{Week 2 -- Foundations: Ethics, law \& responsible collection}
\textcolor{CUGold}{\textbf{August 24, 26 \& 28}}\\
The legal and ethical contours of web data retrieval: terms of service, copyright, privacy, and research ethics. Textbook: \href{\bookurl/blob/main/ch-02-ethics.qmd}{Chapter 2 (Ethics, Law, and Responsible Data Collection)}.

\section{Week 3 -- Foundations: The post-API age}
\textcolor{CUGold}{\textbf{August 31, September 2 \& 4}}\\
Changing contexts of data access, platform transparency, and the ``post-API'' turn. Textbook: \href{\bookurl/blob/main/ch-03-post-api.qmd}{Chapter 3 (The Post-API Age)}.

\section{Week 4 -- Foundations: Data formats (XML \& JSON)}
\textcolor{CUGold}{\textbf{September 9 \& 11}}\\
JSON and XML data formats and supporting libraries. \textit{No class Monday, September 7 (Labor Day); this week the concept and notebook are introduced on Wednesday and we hold the revision workshop on Friday.} Textbook: \href{\bookurl/blob/main/ch-04-data-formats.qmd}{Chapter 4 (Data Formats: XML and JSON)}.

\section{Week 5 -- Foundations: Web architecture \& protocols}
\textcolor{CUGold}{\textbf{September 14, 16 \& 18}}\\
Web protocols like IP and HTTP using libraries like \href{https://docs.python.org/3/library/urllib.html}{\texttt{urllib}} and \href{https://requests.readthedocs.io/}{\texttt{requests}}. Textbook: \href{\bookurl/blob/main/ch-05-protocols.qmd}{Chapter 5 (Web Architecture and Protocols)}.

\section{Week 6 -- Documents: Parsing static web pages}
\textcolor{CUGold}{\textbf{September 21, 23 \& 25}}\\
HTML elements and document parsing on static web pages with \href{https://www.crummy.com/software/BeautifulSoup/bs4/doc/}{\texttt{BeautifulSoup}}. Textbook: \href{\bookurl/blob/main/ch-06-static-pages.qmd}{Chapter 6 (Parsing Static Web Pages)}.

\section{Week 7 -- Documents: Archived web pages}
\textcolor{CUGold}{\textbf{September 28, 30 \& October 2}}\\
Using the Internet Archive's Wayback Machine to analyze archived web pages. Textbook: \href{\bookurl/blob/main/ch-07-archives.qmd}{Chapter 7 (Archived Web Pages and the Wayback Machine)}.

\section{Week 8 -- Documents: Dynamic web pages}
\textcolor{CUGold}{\textbf{October 5, 7 \& 9}}\\
Rendering and scraping dynamic web pages with \href{https://selenium-python.readthedocs.io/}{\texttt{selenium}}. Textbook: \href{\bookurl/blob/main/ch-08-dynamic-pages.qmd}{Chapter 8 (Dynamic Web Pages with Selenium)}.

\section{Week 9 -- Documents: Extracting data from PDFs}
\textcolor{CUGold}{\textbf{October 12, 14 \& 16}}\\
Accessing and parsing PDFs with \href{https://pypdf.readthedocs.io/}{\texttt{pypdf}}. Textbook: \href{\bookurl/blob/main/ch-09-pdfs.qmd}{Chapter 9 (Extracting Data from PDFs)}.

\section{Week 10 -- APIs: Wikipedia}
\textcolor{CUGold}{\textbf{October 19, 21 \& 23}}\\
Retrieving data from Wikipedia's APIs. \textbf{Final Project proposal due Friday, October 23.} Textbook: \href{\bookurl/blob/main/ch-10-wikipedia.qmd}{Chapter 10 (Introduction to APIs: Wikipedia)}.

\section{Week 11 -- APIs: Government data}
\textcolor{CUGold}{\textbf{October 26, 28 \& 30}}\\
Retrieving and filtering data from government APIs like the Census and Federal Reserve. Textbook: \href{\bookurl/blob/main/ch-11-government.qmd}{Chapter 11 (Government Data APIs)}.

\section{Week 12 -- APIs: Social \& media platforms}
\textcolor{CUGold}{\textbf{November 2, 4 \& 6}}\\
Retrieving content, social graphs, and activity streams from social and media platform APIs. Textbook: \href{\bookurl/blob/main/ch-12-social.qmd}{Chapter 12 (Social and Media Platform APIs)}.

\section{Week 13 -- APIs: AI \& language models}
\textcolor{CUGold}{\textbf{November 9, 11 \& 13}}\\
Accessing and using AI and large language models through their APIs. Textbook: \href{\bookurl/blob/main/ch-13-ai-apis.qmd}{Chapter 13 (AI and Language Model APIs)}.

\section{Week 14 -- Practice: Automating data collection}
\textcolor{CUGold}{\textbf{November 16, 18 \& 20}}\\
Automating and scheduling web data retrieval with tools like \href{https://docs.github.com/en/actions}{GitHub Actions}. Textbook: \href{\bookurl/blob/main/ch-14-automation.qmd}{Chapter 14 (Automating Data Collection)}.

\section{Week 15 -- Fall Break}
\textcolor{CUGold}{\textbf{November 23--27}}\\
No class (Fall Break; university closed November 26 \& 27). Use this week to make progress on your Final Project.

\section{Week 16 -- Practice: Research design \& final presentations}
\textcolor{CUGold}{\textbf{November 30, December 2 \& 4}}\\
Matching web data sources with research designs; workshopping and presenting Final Projects. \textit{December 4 is the last day of class and follows a Monday schedule.} Textbook: \href{\bookurl/blob/main/ch-15-research-design.qmd}{Chapter 15 (Research Design with Web Data)}. \textbf{Final Project repository and write-up due Monday, December 7.}

% \newpage


\renewcommand{\bibsection}{\section{\huge \bibname}\prebibhook}
\baselineskip 14.2pt
\nobibliography{refs}
\bibliographystyle{apalike}

\end{document}
